%=========== ΛΥΣΗ - 15ο ΘΕΜΑ Δ ===========
\begin{thema}{Δ}
\begin{erwthma}
\item Η $f$ είναι παραγωγίσιμη στο $\mathbb{R}$ και για κάθε $x\in\mathbb{R}$ είναι
\begin{align*}
f'(x)
&=(x^2e^{-x})'\\
&=2xe^{-x}-x^2e^{-x}\\
&=xe^{-x}(2-x).
\end{align*}
Επειδή $e^{-x}>0$, το πρόσημο της $f'$ είναι το πρόσημο της παράστασης $x(2-x)$. Επομένως:
\begin{itemize}
\item $f'(x)<0$ για $x<0$ ή $x>2$,
\item $f'(x)=0$ για $x=0$ ή $x=2$,
\item $f'(x)>0$ για $0<x<2$.
\end{itemize}
Στον παρακάτω πίνακα βλέπουμε το πρόσημο της $f'$ και τη μονοτονία της $f$:
\begin{center}
\begin{tikzpicture}
\tkzTabInit[color,lgt=2,espcl=1.5,colorC = maincolor!50!white,colorL = maincolor!20!white,colorV = maincolor!50!white]%
{$x$ / .8,$f'(x)$ /0.8,$f(x)$/1}%
{$-\infty$,$0$,$2$,$+\infty$}%
\tkzTabLine{, -, z, +, z, -, }
\tkzTabVar{+/{+\infty},-/{0},+/{\frac{4}{e^2}},-/{0}}
\end{tikzpicture}
\end{center}

Άρα η $f$ είναι \textbf{γνησίως φθίνουσα} στα διαστήματα $(-\infty,0]$ και $[2,+\infty)$ και \textbf{γνησίως αύξουσα} στο $[0,2]$.

Επομένως παρουσιάζει:
\begin{itemize}
\item τοπικό ελάχιστο στο $x=0$, με τιμή $f(0)=0$,
\item τοπικό μέγιστο στο $x=2$, με τιμή
\[
f(2)=4e^{-2}=\frac{4}{e^2}.
\]
\end{itemize}

Για τις ασύμπτωτες της $C_f$ έχουμε
\[
\lim_{x\to+\infty}x^2e^{-x}
=\lim_{x\to+\infty}\frac{x^2}{e^x}
\xlongequal[\text{\eng{DLH}}]{\frac{+\infty}{+\infty}}
\lim_{x\to+\infty}\frac{2x}{e^x}
\xlongequal[\text{\eng{DLH}}]{\frac{+\infty}{+\infty}}
\lim_{x\to+\infty}\frac{2}{e^x}
=0.
\]
Άρα η ευθεία
\[
{y=0}
\]
είναι οριζόντια ασύμπτωτη της $C_f$ στο $+\infty$.

Επίσης
\[
\lim_{x\to-\infty}x^2e^{-x}=+\infty,
\]
ενώ
\[
\lim_{x\to-\infty}\frac{f(x)}{x}
=\lim_{x\to-\infty}xe^{-x}
=-\infty.
\]
Άρα η $C_f$ δεν έχει οριζόντια ή πλάγια ασύμπτωτη στο $-\infty$.

\item Για κάθε $x\in\mathbb{R}$ είναι
\begin{align*}
f''(x)
&=\left[xe^{-x}(2-x)\right]'\\
&=e^{-x}(2-x)-xe^{-x}(2-x)-xe^{-x}\\
&=e^{-x}(x^2-4x+2).
\end{align*}
Επειδή $e^{-x}>0$, το πρόσημο της $f''$ είναι το πρόσημο του τριωνύμου
\[
x^2-4x+2.
\]
Οι ρίζες του είναι
\[
x_{1,2}=\frac{4\pm\sqrt{16-8}}{2}=2\pm\sqrt{2}.
\]
Επομένως:
\begin{itemize}
\item $f''(x)>0$ για $x<2-\sqrt{2}$ ή $x>2+\sqrt{2}$,
\item $f''(x)<0$ για $2-\sqrt{2}<x<2+\sqrt{2}$.
\end{itemize}
Άρα η $f$ είναι κυρτή στα διαστήματα
\[
(-\infty,2-\sqrt{2}]
\quad\text{και}\quad
[2+\sqrt{2},+\infty)
\]
και κοίλη στο
\[
[2-\sqrt{2},2+\sqrt{2}].
\]

Η $f''$ αλλάζει πρόσημο στα $x=2\pm\sqrt{2}$, οπότε η $C_f$ έχει σημεία καμπής τα
\[
{K_1\left(2-\sqrt{2},(2-\sqrt{2})^2e^{-(2-\sqrt{2})}\right)}
\]
και
\[
{K_2\left(2+\sqrt{2},(2+\sqrt{2})^2e^{-(2+\sqrt{2})}\right)}.
\]

\item Στο διάστημα $[0,+\infty)$, από τη μελέτη της μονοτονίας, η $f$ παρουσιάζει ολικό μέγιστο στο $x=2$, με τιμή $\frac{4}{e^2}$. Άρα, για κάθε $x\geq0$,
\begin{align*}
f(x)&\leq f(2)\\
\Leftrightarrow x^2e^{-x}&\leq4e^{-2}\\
\Leftrightarrow {x^2\leq4e^{x-2}}.
\end{align*}
Η ισότητα ισχύει μόνο για $x=2$.

\item Για $a>0$ έχουμε
\[
E(a)=\int_0^a x^2e^{-x}\d x.
\]
Με διαδοχική ολοκλήρωση κατά παράγοντες βρίσκουμε
\begin{align*}
\int x^2e^{-x}\d x
&=-x^2e^{-x}+2\int xe^{-x}\d x\\
&=-x^2e^{-x}-2xe^{-x}-2e^{-x}+c\\
&=-(x^2+2x+2)e^{-x}+c.
\end{align*}
Επομένως
\begin{align*}
E(a)
&=\left[-(x^2+2x+2)e^{-x}\right]_0^a\\
&=2-(a^2+2a+2)e^{-a}.
\end{align*}
Άρα
\begin{align*}
\lim_{a\to+\infty}E(a)
&=2-\lim_{a\to+\infty}\frac{a^2+2a+2}{e^a}\\
&={2},
\end{align*}
αφού η εκθετική συνάρτηση υπερισχύει κάθε πολυωνυμικής στο $+\infty$.
\end{erwthma}
\end{thema}
%=========== ΤΕΛΟΣ ΛΥΣΗΣ - 15ο ΘΕΜΑ Δ ===========
